\documentclass{ximera}

\input{../../preamble.tex}

\title{Trigonometric}

\begin{document}

\begin{abstract}
oscillation
\end{abstract}
\maketitle





\subsection*{The Core Sine}


We get the Core sine function by viewing the $y$-coordinates of the unit circle as a function of the angle.




\begin{image}
\begin{tikzpicture} 
  \begin{axis}[
            domain=-1.5:1.5, ymax=1.5, xmax=1.5, ymin=-1.5, xmin=-1.5,
            axis lines =center, unit vector ratio*=1 1 1, xlabel=$x$, ylabel=$y$,
            ytick={-10,-8,-6,-4,-2,2,4,6,8,10},
            xtick={-10,-8,-6,-4,-2,2,4,6,8,10},
            ticklabel style={font=\scriptsize},
            every axis y label/.style={at=(current axis.above origin),anchor=south},
            every axis x label/.style={at=(current axis.right of origin),anchor=west},
            axis on top
          ]
          
            \addplot [line width=2, penColor, smooth,samples=100,domain=(0:6.3)] ({cos(deg(x))},{sin(deg(x)});


          \addplot [color=penColor2,only marks,mark=*] coordinates{(0.5,0.866)};


          %\draw[decoration={brace,raise=.2cm,mirror},decorate,thin] (axis cs:0.75,0)--(axis cs:0.75,0.866);
          %\draw[decoration={brace,raise=.2cm},decorate,thin] (axis cs:0,0.95)--(axis cs:0.75,0.95);
          %\node[anchor=east] at (axis cs:1.85,3) {$d$};
          %\node[anchor=east] at (axis cs:4.6,1) {$f(d)$};
                     
          \node at (axis cs:0.75,1.2) [penColor2] {$(\cos(\theta),\sin(\theta))$};


          \addplot [line width=1, penColor2, smooth,samples=100,domain=(0:0.5)] ({x},{1.7320*x});
          \node at (axis cs:0.25,0.15) [penColor2] {$\theta$};


          %\addplot [color=penColor,only marks,mark=*] coordinates{(-4,-1) (0,1) (1,-6.5) (7,-3.5)};

          %\addplot [line width=1, penColor2, smooth,samples=100,domain=(-4:0)] ({x},{0});
          %\addplot [line width=1, penColor2, smooth,samples=100,domain=(1:7)] ({x},{0});
           

  \end{axis}
\end{tikzpicture}
\end{image}




As you travel counterclockwise around the circle from the point $(1,0)$, the $y$-coordinate of the points

\begin{itemize}
\item increase from $0$ to $1$, then
\item decrease from $1$ to $0$ and then to $-1$, then
\item increase from $-1$ to $0$
\end{itemize}

This keeps repeating with each full circle, counterclockwise (positive angles) or clockwise (negative angles).




\begin{image}
\begin{tikzpicture}
  \begin{axis}[
            domain=-10:10, ymax=1.5, xmax=10, ymin=-1.5, xmin=-10,
            axis lines =center, xlabel={$\theta$}, ylabel={$y$}, grid = major, grid style={dashed},
            ytick={-1.5,-1,-0.5,0.5,1,1.5},
            xtick={-7.85, -6.28, -4.71, -3.14, -1.57, 0, 1.57, 3.142, 4.71, 6.28, 7.85},
            xticklabels={$\tfrac{-5\pi}{2}$,$-2\pi$,$\tfrac{-3\pi}{2}$,$-\pi$, $\tfrac{-\pi}{2}$, $0$, $\tfrac{\pi}{2}$, $\pi$, $\tfrac{3\pi}{2}$, $2\pi$, $\tfrac{5\pi}{2}$},
            yticklabels={$1.5$,$-1$,$-0.5$,$0.5$,$1$,$1.5$}, 
            ticklabel style={font=\scriptsize},
            every axis y label/.style={at=(current axis.above origin),anchor=south},
            every axis x label/.style={at=(current axis.right of origin),anchor=west},
            axis on top
          ]
          

            \addplot [line width=2, penColor, smooth,samples=300,domain=(-10:10),<->] {sin(deg(x)};
            %\addplot [line width=2, penColor2, smooth,samples=300,domain=(-10:10),<->] {cos(deg(x)};


            \node[penColor] at (axis cs:3.5,1.1) [anchor=north] {$\sin(\theta)$};
            %\node[penColor2] at (axis cs:1.4,1.3) [anchor=north] {$\cos(\theta)$};



  \end{axis}
\end{tikzpicture}
\end{image}



The general sine function will be seen as compositions of the Core sine function and two linear functions.





\begin{summary} \textbf{\textcolor{red!70!yellow}{Sine Analysis}}


\textbf{Domain:} The domain of the Core sine function is all real numbers, $(-\infty, \infty)$. 

\textbf{Zeros:} The zeros of the Core sine function are $\{ k\pi \, | \, k \in \mathbb{Z} \}$.

\textbf{Continuity:} The Core sine function is continuous.

\textbf{End-Behavior:} The Core sine function keeps oscillating between $1$ and $-1$.



\textbf{Behavior (Increasing and Decreasing):} The Core sine function switches between increasing and decreasing with a period of $2\pi$.

On the interval $[0 , 2\pi)$, the Core sine function

\begin{itemize}
\item increases on $\left(0, \frac{\pi}{2} \right)$
\item increases on $\left(\frac{\pi}{2},  \frac{3\pi}{2}\right)$
\item increases on $\left(\frac{3\pi}{2}, 2\pi \right)$
\end{itemize}








\textbf{Global Maximum and Minimum:} 

The global maximum value is $1$, which occurs at $\{ \frac{\pi}{2} + 2k\pi \, | \, k \in \mathbb{Z} \}$.


The global minimum value is $-1$, which occurs at $\{ \frac{3\pi}{2} + 2k\pi \, | \, k \in \mathbb{Z} \}$.




\textbf{Local Maximums and Minimums:}  The global maximum and minimum are the only local extrema.




\textbf{Range:} The range of the Core sine function is $[-1, 1]$.


\end{summary}
























\subsection*{The Core Cosine}


We get the Core cosine function by viewing the $x$-coordinates of the unit circle as a function of the angle.



As you travel counterclockwise around the circle from the point $(1,0)$, the $x$-coordinate of the points

\begin{itemize}
\item decrease from $1$ to $0$ and then to $-1$, then
\item increase from $-1$ to $0$ and then to $1$
\end{itemize}

This keeps repeating with each full circle, counterclockwise (positive angles) or clockwise (negative angles).




\begin{image}
\begin{tikzpicture}
  \begin{axis}[
            domain=-10:10, ymax=1.5, xmax=10, ymin=-1.5, xmin=-10,
            axis lines =center, xlabel={$\theta$}, ylabel={$y$}, grid = major, grid style={dashed},
            ytick={-1.5,-1,-0.5,0.5,1,1.5},
            xtick={-7.85, -6.28, -4.71, -3.14, -1.57, 0, 1.57, 3.142, 4.71, 6.28, 7.85},
            xticklabels={$\tfrac{-5\pi}{2}$,$-2\pi$,$\tfrac{-3\pi}{2}$,$-\pi$, $\tfrac{-\pi}{2}$, $0$, $\tfrac{\pi}{2}$, $\pi$, $\tfrac{3\pi}{2}$, $2\pi$, $\tfrac{5\pi}{2}$},
            yticklabels={$1.5$,$-1$,$-0.5$,$0.5$,$1$,$1.5$}, 
            ticklabel style={font=\scriptsize},
            every axis y label/.style={at=(current axis.above origin),anchor=south},
            every axis x label/.style={at=(current axis.right of origin),anchor=west},
            axis on top
          ]
          

            %\addplot [line width=2, penColor, smooth,samples=300,domain=(-10:10),<->] {sin(deg(x)};
            \addplot [line width=2, penColor2, smooth,samples=300,domain=(-10:10),<->] {cos(deg(x)};


            %\node[penColor] at (axis cs:3.5,1.1) [anchor=north] {$\sin(\theta)$};
            \node[penColor2] at (axis cs:1.4,1.3) [anchor=north] {$\cos(\theta)$};



  \end{axis}
\end{tikzpicture}
\end{image}



The general cosine function will be seen as compositions of the Core cosine function and two linear functions.





\begin{summary} \textbf{\textcolor{red!70!yellow}{Cosine Analysis}}


\textbf{Domain:} The domain of the Core cosine function is all real numbers, $(-\infty, \infty)$. 

\textbf{Zeros:} The zeros of the Core sine function are $\{ \frac{\pi}{2}  + k\pi \, | \, k \in \mathbb{Z} \}$.

\textbf{Continuity:} The Core cosine function is continuous.

\textbf{End-Behavior:} The Core cosine function keeps oscillating between $1$ and $-1$.



\textbf{Behavior (Increasing and Decreasing):} The Core cosine function switches between increasing and decreasing with a period of $2\pi$.

On the interval $[0 , 2\pi)$, the Core cosine function

\begin{itemize}
\item decreases on $(0, \pi)$
\item increases on $(\pi, 2\pi)$
\end{itemize}








\textbf{Global Maximum and Minimum:} 

The global maximum value is $1$, which occurs at $\{ 2k\pi \, | \, k \in \mathbb{Z} \}$.


The global minimum value is $-1$, which occurs at $\{  (2k+1)\pi \, | \, k \in \mathbb{Z} \}$.




\textbf{Local Maximums and Minimums:}  The global maximum and minimum are the only local extrema.




\textbf{Range:} The range of the Core cosine function is $[-1, 1]$.


\end{summary}



























\subsection*{The Core Tangent}


Rational functions were quotients of polynomials.  We can also create quotients with trigonometric functions.


\begin{definition}  \textbf{\textcolor{green!50!black}{Tangent}} 


\textbf{Tangent} is the quotient of sine and cosine


\[   \tan(\theta) = \frac{\sin(\theta)}{\cos(\theta)}      \]



Its domain is all real numbers except the zeros of cosine, $\left\{  \frac{(2k+1)\pi}{2}   \, | \,   k \in \textbf{Z}   \right\}$.

\end{definition}


At each of these singularities, the graph has a vertical asymptote, where the function becomes unbounded.




Graph of $z = \tan(\theta)$.





\begin{image}
\begin{tikzpicture}
  \begin{axis}[
            xmin=-6.75,xmax=6.75,ymin=-5,ymax=5,
            axis lines=center,
            width=6in,
            height=3in,
            xtick={-6.28, -4.71, -3.14, -1.57, 0, 1.57, 3.142, 4.71, 6.28},
            xticklabels={$-2\pi$,$\frac{-3\pi}{2}$,$-\pi$, $\frac{-\pi}{2}$, $0$, $\frac{\pi}{2}$, $\pi$, $\frac{3\pi}{2}$, $2\pi$},       
            unit vector ratio*=1 1 1,
            xlabel=$\theta$, ylabel=$z$,
            ticklabel style={font=\scriptsize},
            every axis y label/.style={at=(current axis.above origin),anchor=south},
            every axis x label/.style={at=(current axis.right of origin),anchor=west},
          ]        
          \addplot [line width=2, penColor, samples=200,smooth, domain=(-1.37:1.37),<->] {tan(deg(x))};
          \addplot [line width=2, penColor, samples=100,smooth, domain=(-4.51:-1.76),<->] {tan(deg(x))};
          \addplot [line width=2, penColor, samples=100,smooth, domain=(-6.08:-4.9),<->] {tan(deg(x))};
          \addplot [line width=2, penColor, samples=100,smooth, domain=(1.76:4.51),<->] {tan(deg(x))};
          \addplot [line width=2, penColor, samples=100,smooth, domain=(4.9:6.75),<->] {tan(deg(x))};
          
          \addplot [textColor,dashed] plot coordinates {(-4.71,-5) (-4.71,5)};
          \addplot [textColor,dashed] plot coordinates {(-1.57,-5) (-1.57,5)};
          \addplot [textColor,dashed] plot coordinates {(1.57,-5) (1.57,5)};
          \addplot [textColor,dashed] plot coordinates {(4.71,-5) (4.71,5)};
          
          %\node at (axis cs:.4,1.25) [penColor3] {$\tan(\theta)$};    
        \end{axis}
\end{tikzpicture}

\end{image}



Tangent takes on every value in $\mathbb{R}$ exactly once inside the open interval $\left( \frac{-\pi}{2}, \frac{\pi}{2} \right)$.  It is then periodic with a period of $\pi$.

Tangent is always \wordChoice{\choice[correct]{increasing} \choice{decreasing}} inside its domain.

Tangent is unbounded and has no global or local maximums or minimums.

Tangent has an infinite number of zeros. They coincide with the zeros of \wordChoice{\choice[correct]{sine} \choice{cosine}} : $\{  k \pi    \, | \,   k \in \textbf{Z}        \}$















\begin{summary} \textbf{\textcolor{red!70!yellow}{Tangent Analysis}}


\textbf{Domain:} The domain of the Core tangent function is all real, except the zeros of cosine. 

\[ (-\infty, \infty) -  \{ \frac{\pi}{2}  + k\pi \, | \, k \in \mathbb{Z} \} \]


Another way to say this is the domain is the union of all intervals of the form 

\[ \left(  -\frac{\pi}{2}  + k\pi, \frac{\pi}{2}  + k\pi   \right) \, | \, k \in \mathbb{Z}   \]


These real numbers left out of the domain are singularities of tangent.



\textbf{Zeros:} The zeros of the Core tangent are the zeros of the Core sine function, which are $\{ \frac{\pi}{2}  + k\pi \, | \, k \in \mathbb{Z} \}$.

\textbf{Continuity:} The Core tangent function is continuous.

\textbf{End-Behavior:} The Core tangent function keeps repeating with a period of $\pi$.

It does become unbounded near its singularities.


\[
\lim\limits_{\theta \to -\frac{\pi}{2}^+} \tan(\theta) = -\infty
\]


\[
\lim\limits_{\theta \to \frac{\pi}{2}^-} \tan(\theta) = \infty
\]





\textbf{Behavior (Increasing and Decreasing):} The Core tangent function is an increasing function.





\textbf{Global Maximum and Minimum:} The Core tangent function has no global maximum or minimum.




\textbf{Local Maximums and Minimums:}  The Core tangent function has no local maximums or minimums.




\textbf{Range:} The range of the Core cosine function is $[-1, 1]$.


\end{summary}
































\begin{center}
\textbf{\textcolor{green!50!black}{ooooo-=-=-=-ooOoo-=-=-=-ooooo}} \\

more examples can be found by following this link\\ \link[More Examples of Analysis]{https://ximera.osu.edu/csccmathematics/precalculus1/precalculus1/libraryAnalysis2/examples/exampleList}

\end{center}




\end{document}
